% 除法定则（综述）
% license CCBYSA3
% type Wiki

本文根据 CC-BY-SA 协议转载翻译自维基百科\href{https://en.wikipedia.org/wiki/Quotient_rule}{相关文章}

在微积分中，商法则是一种求导方法，用于求两个可微函数之比的导数。设$h(x) = \frac{f(x)}{g(x)}$，其中$f$和$g$都是可微函数，且$g(x) \neq 0$。商法则表明，$h(x)$ 的导数为：
$$
h'(x) = \frac{f'(x)g(x) - f(x)g'(x)}{(g(x))^2}.~
$$
这一法则可以通过使用其他求导规则以多种方式证明。
\subsection{例子}
\subsubsection{例 1：基础例子}
已知$h(x) = \frac{e^x}{x^2},$令$f(x) = e^x, \quad g(x) = x^2$，那么根据商法则：
$$
\begin{aligned}
\frac{d}{dx}\left(\frac{e^x}{x^2}\right) 
&= \frac{\left(\frac{d}{dx} e^x\right)(x^2) - (e^x)\left(\frac{d}{dx} x^2\right)}{(x^2)^2}\\
&= \frac{(e^x)(x^2) - (e^x)(2x)}{x^4}\\
&= \frac{x^2 e^x - 2x e^x}{x^4}\\
&= \frac{x e^x - 2 e^x}{x^3}\\
&= \frac{e^x (x-2)}{x^3}.\\
\end{aligned}~
$$
\subsubsection{例 2：正切函数的导数}
商法则可以用来求$\tan x$的导数。由于$\tan x = \frac{\sin x}{\cos x}$,

所以：
$$
\begin{aligned}
\frac{d}{dx}\tan x &= \frac{d}{dx}\left(\frac{\sin x}{\cos x}\right)\\
&= \frac{\left(\frac{d}{dx}\sin x\right)(\cos x) - (\sin x)\left(\frac{d}{dx}\cos x\right)}{\cos^2 x}\\
&= \frac{(\cos x)(\cos x) - (\sin x)(- \sin x)}{\cos^2 x}\\
&= \frac{\cos^2 x + \sin^2 x}{\cos^2 x}\\
&= \frac{1}{\cos^2 x} = \sec^2 x.\\
\end{aligned}~
$$
\subsection{倒数法则}
倒数法则是商法则的一个特殊情形，此时分子为$f(x) = 1$,应用商法则可得：
$$
h'(x) = \frac{d}{dx}\left[\frac{1}{g(x)}\right] 
= \frac{0 \cdot g(x) - 1 \cdot g'(x)}{(g(x))^2} 
= \frac{-g'(x)}{(g(x))^2}.~
$$
利用链式法则也能得到相同的结果。

